\chapter{Einfluss künstlicher Intelligenz auf Projektpraktiken}
\label{sec:KIProjektpraktiken}

\section{Aktuelle Einsatzfelder künstlicher Intelligenz}
\label{sec:AktuelleKIEinsatzfelder}

\subsection{Projektplanung und Aufwandsschätzung}
\subsection{Risikoerkennung und Entscheidungsunterstützung}
\subsection{Ressourcen-, Termin- und Kostenmanagement}
\subsection{Kommunikation, Dokumentation und Wissensmanagement}
\subsection{Qualitätssicherung und Projektcontrolling}

\section{Einfluss auf klassische, agile und hybride Vorgehensmodelle}
\label{sec:KIEinflussVorgehensmodelle}

\subsection{Veränderungen klassischer Projektpraktiken}
\subsection{Veränderungen agiler Projektpraktiken}
\subsection{Potenziale für hybride Projektorganisationen}

\section{Veränderung von Rollen und Kompetenzen}
\label{sec:RollenUndKompetenzen}

\subsection{Projektleitung und Projektteam}
\subsection{Mensch-KI-Zusammenarbeit}
\subsection{Erforderliche fachliche und überfachliche Kompetenzen}

\section{Chancen und Nutzenpotenziale}
\label{sec:KIChancen}

\section{Risiken und Herausforderungen}
\label{sec:KIRisiken}

\subsection{Datenqualität, Nachvollziehbarkeit und Fehlentscheidungen}
\subsection{Datenschutz, Informationssicherheit und Compliance}
\subsection{Akzeptanz, Verantwortung und ethische Fragestellungen}

\section{Zukünftige Projektpraktiken}
\label{sec:ZukuenftigeProjektpraktiken}

\subsection{Automatisierung und autonome Projektassistenz}
\subsection{Prädiktive und adaptive Projektsteuerung}
\subsection{Zukunftsszenarien für das IT-Projektmanagement}
